%======================================================================
% T9_mg_cooling.pdf
%======================================================================
\section{Cvičení: Aplikace Maxwellových relací -- magnetické chlazení}
\strana{T9mg-1}

\begin{priklad}[Příklad 1 -- Magnetické chlazení]
Paramagnetická látka $\big[M(H=0) = 0,\ \dd W = H\,\dd M\big]$ je charakterizována kalorickou rovnicí $U = BT^4$
a izotermickou susceptibilitou $\chi_T = \dfrac{A}{T}$, $A > 0$. \quad ($\chi_T$ $\updownarrow$ id.\ paramagnet)
\begin{enumerate}[label=\arabic*)]
  \item Kolik tepla přiteče do systému, pokud izotermicky zapneme pole $H = 0 \to H = H_0$?
  \item Jak se změní teplota, pokud následně adiabaticky mg.\ pole vypneme?
\end{enumerate}
\end{priklad}

\underline{Pozn:}
\begin{itemize}[nosep]
  \item $T = \text{konst} \Leftrightarrow$ pomalé zapínání (vibrační a mg.\ stupně volnosti stihají interagovat, tep.\ interakce
        s okolím dostatečně rychlá k udržení konst.\ teploty)
  \item $\dd Q = 0 \Leftrightarrow$ rychlé vypnutí (teplo přitéct nestihne)
\end{itemize}

\noindent\rule{\linewidth}{0.4pt}

\subsection*{Řešení}
{\color{NavyBlue}
1) \ $U = BT^4 \Rightarrow \dd T = 0 \Leftrightarrow \dd U = 0 \Rightarrow Q = -W$
\[
  Q = -\int_0^{H_0}H\,\dd M = -\int_0^{H_0}H\Big(\frac{\partial M}{\partial H}\Big)_T\dd H = -\frac{A}{T}\int_0^{H_0}H\,\dd H = -\frac{AH_0^2}{T}
\]
\quad $\Rightarrow$ zapnutím pole jsme teplo ze systému odčerpali (orientace spinů $\Rightarrow$ pokles entropie)
\pozn{v originále takto; $\int_0^{H_0}H\,\dd H = H_0^2/2$}
}

%----------------------------------------------------------------------
\strana{T9mg-2}
{\color{NavyBlue}
2) \
\[
  \Delta T = \int_{H_0}^{0}\Big(\frac{\partial T}{\partial H}\Big)_S\dd H
  = \int_0^{H_0}\frac{\big(\frac{\partial S}{\partial H}\big)_T}{\big(\frac{\partial S}{\partial T}\big)_H}\,\dd H
\]
\quad $\Rightarrow$ potřebujeme nalézt $S = S(T,H)$ \\
\quad $\Rightarrow$ potom ale lze $T_f$ určit ze vztahu \quad $S(T_f, 0) = S(T_i, H_0)$

a) \ FR je $U = U(S,M)$, přitom známe \quad $U(T,M) = U\big(S(T,M),M\big) = BT^4$
\[
  \Rightarrow\ \Big(\frac{\partial U}{\partial M}\Big)_T = 0
  = \Big(\frac{\partial U}{\partial S}\Big)_M\Big(\frac{\partial S}{\partial M}\Big)_T + \Big(\frac{\partial U}{\partial M}\Big)_S
  = T\Big(\frac{\partial S}{\partial M}\Big)_T + H
  \qquad\Rightarrow\ \boxed{\Big(\frac{\partial S}{\partial M}\Big)_T = -\frac{H}{T}}
\]

b) \ $S(T,M) = S\big(T, H(T,M)\big)$
\[
  \Rightarrow\ \Big(\frac{\partial S}{\partial M}\Big)_T = \Big(\frac{\partial S}{\partial H}\Big)_T\Big(\frac{\partial H}{\partial M}\Big)_T
  = \Big(\frac{\partial S}{\partial H}\Big)_T\Big(\frac{A}{T}\Big)^{-1}
  \qquad\Rightarrow\ \vzorec{\Big(\frac{\partial S}{\partial H}\Big)_T = -\frac{AH}{T^2}}\ \checkmark
\]

c) \ $\Big(\dfrac{\partial S}{\partial T}\Big)_H = \dfrac{C_H}{T}$ \quad \& \quad
$C_M = \Big(\dfrac{\partial U}{\partial T}\Big)_M = 4BT^3$ \\
\quad $\Rightarrow$ potřebuji Meyerův vztah
\begin{align*}
  C_H = T\Big(\frac{\partial S}{\partial T}\Big)_H &= \Big/\,S(T,H) = S\big(T, M(T,H)\big)\Big/ \\
  &= T\Big(\frac{\partial S}{\partial T}\Big)_M + T\Big(\frac{\partial S}{\partial M}\Big)_T\Big(\frac{\partial M}{\partial T}\Big)_H
   = \Big/\,\dd F = -S\,\dd T + H\,\dd M\Big/
\end{align*}
}

%----------------------------------------------------------------------
\strana{T9mg-3}
{\color{NavyBlue}
\[
  = C_M - T\Big(\frac{\partial H}{\partial T}\Big)_M\Big(\frac{\partial M}{\partial T}\Big)_H
  = C_M + \frac{\big(\frac{\partial M}{\partial T}\big)_H^2}{\big(\frac{\partial M}{\partial H}\big)_T}
  = C_M + \frac{T\big(\frac{\partial M}{\partial T}\big)_H^2}{\chi_T}
\]
\begin{itemize}[nosep]
  \item $\chi_T = \dfrac{A}{T}$; \quad $C_M = 4BT^3$
  \item $\Big(\dfrac{\partial M}{\partial T}\Big)_H = \Big/\,\dd G = -S\,\dd T - M\,\dd H\Big/ = \Big(\dfrac{\partial S}{\partial H}\Big)_T = -\dfrac{AH}{T^2}$ \\
        $\to$ samozřejmě není náhoda, že se \uv{kruh uzavírá}, neboť $\big(\frac{\partial S}{\partial H}\big)_T$ a
        $\big(\frac{\partial S}{\partial T}\big)_H$ nejsou zcela nezávislé -- $\dd S(H,T)$ je úplný diferenciál
\end{itemize}

$\Rightarrow$ i) \ $\Big(\dfrac{\partial S}{\partial T}\Big)_H = 4BT^2 + \dfrac{A^2H^2}{T^4}\,\dfrac{T}{A} = 4BT^2 + \dfrac{AH^2}{T^3}$

\phantom{$\Rightarrow$} ii) \ $\Big(\dfrac{\partial S}{\partial H}\Big)_T = -\dfrac{AH}{T^2}$
\hfill $\Big(\dfrac{\partial^2S}{\partial T\partial H}\Big) = \dfrac{2AH}{T^3} \overset{?}{=} \Big(\dfrac{\partial^2S}{\partial H\partial T}\Big) = \dfrac{2AH}{T^3}\ \checkmark$

\phantom{$\Rightarrow$} ii) \ $S = -\dfrac{AH^2}{2T^2} + f(T)$

\qquad $\Big(\dfrac{\partial S}{\partial T}\Big) = \dfrac{AH^2}{T^3} + f'(T) = \dfrac{AH^2}{T^3} + 4BT^2$

\begin{minipage}{0.55\textwidth}
\[
  \Rightarrow\ \vzorec{S = \frac43BT^3 - \frac{AH^2}{2T^2} + S_0}
\]
\end{minipage}
\begin{minipage}{0.4\textwidth}
$\bullet$ vypadá rozumně: \\
\quad $H\nearrow \Leftrightarrow S\searrow$ \\
\quad $T\nearrow \Leftrightarrow S\nearrow$
\end{minipage}
\[
  \Rightarrow\ \frac43BT_f^3 = \frac43BT_i^3 - \frac{AH_0^2}{2T_i^2}
  \quad\Rightarrow\quad \vzorec{T_f^3 = T_i^3 - \frac{3AH_0^2}{8BT_i^2}\ \Rightarrow\ T_f < T_i}
\]
\begin{itemize}
  \item princip chlazení je, že po vypnutí pole jsou spiny uspořádanější, než by odpovídalo \uv{vibrační} teplotě \\
        $\Rightarrow$ směrem k rovnováze spiny odčerpají teplo z vibrací \\
        (resp.\ v a) teplo odteklo a v b) už se zpět nedostalo; b může probíhat i kvazistat.\ při $\dd S = 0$)
\end{itemize}
}
